\section*{Task 2-1 — Cancelled Orders \& Promotions (Spark DataFrame API)}

\subsection*{(0) Hiểu và bóc đề}
Với mỗi \textbf{thành phố}, tính \% đơn \emph{Cancelled + Standard} thoả \textbf{đồng thời}:
\begin{enumerate}
  \item có $\ge 3$ promotion \textbf{hợp lệ} (active period $\ge 2$ ngày; kể cả mã Amazon);
  \item \texttt{amount} $<$ trung bình đơn \emph{Merchant} + \emph{Courier=Shipped} của \textbf{bang} chứa thành phố đó.
\end{enumerate}
Hai mức địa lý khác nhau trong một câu: ngưỡng tính theo \textbf{bang}, kết quả gộp theo \textbf{thành phố}.
Dataset không có ngày hiệu lực KM — active period phải \emph{suy từ dữ liệu}: mỗi mã, ngày xuất hiện đầu/cuối trên toàn bộ đơn.

\textbf{Dữ liệu đầu vào:} nhóm tiền xử lý \texttt{Amazon Sale Report.csv} bằng \texttt{clean.ipynb} $\to$ \texttt{asr.csv}: chuẩn hoá \texttt{ship-state} (gộp các cách viết trùng như \texttt{NEW DELHI}$\to$\texttt{DELHI}, \texttt{ORISSA}$\to$\texttt{ODISHA}, \texttt{RJ}$\to$\texttt{RAJASTHAN}, sửa lỗi chính tả) và loại 33 dòng thiếu \texttt{ship-state} (128.975 $\to$ 128.942 dòng, 47 $\to$ 37 bang). State normalization làm ở tiền xử lý, không lặp lại trong Spark.

\subsection*{(1) Phân rã thành 4 khối + lý do}
\begin{center}
\begin{tabular}{clll}
\hline
Khối & Việc & Phép & Kết quả \\
\hline
A & Tuổi thọ mã KM & \texttt{groupBy(promo).agg(min,max)}, lọc $\ge 2$ ngày & 284 $\to$ \textbf{185} mã \\
B & Ngưỡng bang & \texttt{Merchant AND Courier=Shipped}, \texttt{groupBy(state).avg} & \textbf{36} bang \\
C & Tập đơn + đếm mã hợp lệ & lọc ứng viên; LEFT join A; LEFT join B & \textbf{6.906} đơn \\
D & Gộp thành phố & \texttt{groupBy(city)}, \% đơn thoả & \textbf{1.434} city \\
\hline
\end{tabular}
\end{center}

\textbf{Lý do thiết kế:}
\begin{itemize}
  \item Khối A dùng \texttt{groupBy().agg(min,max)} \emph{không} window function: cần thu 129k dòng về 284 mã (bảng nhỏ đi), window giữ nguyên số dòng và thêm một shuffle vô ích.
  \item \textbf{LEFT join} (không INNER): đơn không có mã KM vẫn phải nằm ở \emph{mẫu số} của phần trăm; INNER làm mất đơn $\to$ \% sai. Sau join, \texttt{n\_valid} null $\to$ \texttt{coalesce(0)}.
  \item Mã Amazon \textbf{vẫn đếm} (270/284 mã) — đề ghi rõ ``all promotions including those issued by Amazon''.
  \item Chuẩn hoá bang làm ở tiền xử lý (\texttt{clean.ipynb}, 47 $\to$ 37 bang). Task 2-1 chỉ \texttt{upper(trim)} thành phố. Trong tập ứng viên: 1.639 cách viết \texttt{ship-city} thô $\to$ \textbf{1.434} city chuẩn hoá. Tập ứng viên còn \textbf{6.906} đơn (các đơn thiếu \texttt{ship-state} đã bị loại ở tiền xử lý), mọi đơn đều có \texttt{ship-city} $\to$ mẫu số \textbf{6.906}.
  \item Khoá đơn dùng cột \texttt{index} có sẵn (int duy nhất 0..128974), \emph{không} dùng \texttt{monotonically\_increasing\_id} (không ổn định giữa các lần đánh giá).
\end{itemize}

\subsection*{(2) Kết quả = 0\% mọi thành phố — có chứng minh}
Tử số bằng 0 ở mọi thành phố. Kết quả: 1.434 city, mẫu số 6.906 đơn. Đây \emph{là} đáp án đúng, không phải lỗi. Chứng minh bằng cách nới dần điều kiện:
\begin{center}
\begin{tabular}{lrrr}
\hline
Cách hiểu ``cancelled'' & Số đơn & Max mã/đơn & $\ge 3$ mã \\
\hline
Status chứa cancelled + Standard & 6.906 & 0 & \textbf{0} \\
Courier=Cancelled + Standard & 43 & 0 & \textbf{0} \\
Status chứa cancelled, mọi level & 18.332 & 1 & \textbf{0} \\
Courier=Cancelled, mọi level & 5.935 & 2 & \textbf{0} \\
\hline
\end{tabular}
\end{center}
Toàn bộ 6.906 đơn Cancelled+Standard có \texttt{promotion-ids} rỗng; nới hết mức thì max chỉ 2 mã/đơn $\to$ điều kiện $\ge 3$ không bao giờ chạm được. \textbf{Không} hạ ngưỡng để ``ra số đẹp''. File xuất giữ cột chứng minh: \texttt{city, total\_orders, qualified\_orders, pct\_cancelled}.

\textbf{Join có chạy đúng không?} Áp logic đếm mã hợp lệ lên \emph{toàn bộ} dataset (không chỉ đơn cancelled) cho \texttt{max(n\_valid) = 36} — có đơn tới 36 mã hợp lệ. Chứng tỏ join đúng; số 0 ở nhóm cancelled là do \emph{dữ liệu}, không phải lỗi pipeline.

\subsection*{(3) explain(true) — phân tích kế hoạch thực thi}
Chạy pipeline dưới 2 cấu hình (AQE tắt để đọc plan tĩnh; xem cảnh báo cuối mục):

\begin{center}
\begin{tabular}{lcc}
\hline
Đọc từ Physical Plan & Mặc định & \texttt{autoBroadcastJoinThreshold=-1} \\
\hline
(a) Physical join strategy & 3 $\times$ \textbf{BroadcastHashJoin} & 3 $\times$ \textbf{SortMergeJoin} \\
(b) Exchange (shuffle) & 5 & 8 \\
\quad BroadcastExchange & 3 & 0 \\
\quad Sort node & 1 & 7 \\
(c) Số stage của action \texttt{write} (distinct stageIds) & 11 & --- \\
\hline
\end{tabular}
\end{center}

\textbf{(a)} 3 bảng đi join đều tí hon (185 mã hợp lệ, bảng đếm mã/đơn, 36 ngưỡng bang) $<$ 10 MB $\to$ Spark broadcast cả 3, khỏi shuffle bảng đơn. Ép \texttt{threshold=-1} buộc SortMergeJoin: mỗi join thêm 1 Exchange + 2 Sort hai phía (Sort: $1 \to 7$).

\textbf{(b)} Mặc định: 5 shuffle Exchange = 4 \texttt{groupBy} (tuổi thọ mã KM, đếm mã hợp lệ mỗi đơn, ngưỡng bang, kết quả city) + 1 \texttt{rangepartitioning} của \texttt{orderBy} cuối. Broadcast không tạo shuffle bảng lớn nên chỉ có 3 BroadcastExchange bảng nhỏ. Tắt broadcast: 8 Exchange (3 join giờ phải shuffle cả hai phía) và 7 Sort.

\textbf{(c)} Đo bằng \texttt{SparkContext.statusTracker}: đặt action \texttt{write} vào job group \texttt{"task21-write"}, sau đó lấy DISTINCT \texttt{stageIds} của các job trong group đó (\texttt{getJobIdsForGroup} $\to$ \texttt{SparkJobInfo.stageIds}). Kết quả: \textbf{11 stage} (AQE tắt, broadcast bật) — số này khớp kế hoạch tĩnh mặc định và tái lập được mỗi lần chạy (in ra \texttt{writeStages=11}). Việc verify (\texttt{count}/\texttt{sum}/\texttt{max}) đọc lại parquet đã ghi \emph{sau} khi đo, nên không lẫn vào stage count. Với AQE bật, stage có thể coalesce động — đọc thêm ở Spark UI tab Jobs (\texttt{localhost:4040}).

\textbf{Cảnh báo AQE:} Spark 3.5 bật Adaptive Query Execution mặc định $\to$ \texttt{explain(true)} in \emph{kế hoạch tĩnh}; runtime AQE có thể tự đổi (vd re-broadcast). Kế hoạch \emph{thật} xem ở Spark UI tab SQL. Bảng trên chạy với \texttt{spark.sql.adaptive.enabled=false} để plan tĩnh phản ánh đúng khác biệt broadcast.

Plan đầy đủ: \texttt{report/plans/task21\_default.txt}, \texttt{task21\_nobroadcast.txt}.
